\documentclass[10pt,a4paper]{amsart}
\usepackage[margin=19mm]{geometry}
\usepackage{amsmath,amssymb,mathtools}
\usepackage[scr=rsfs,cal=euler]{mathalfa}
\usepackage{xfrac}
\usepackage[unicode,hidelinks,bookmarks=false]{hyperref}
\newcommand{\ie}{i.e.\ }
\newcommand{\cf}{cf.\ }
\newcommand{\resp}{respectively\ }
\newcommand{\Subsection}{\subsection}
\providecommand{\nobreakdash}{\nobreak\hbox{-}}
\setlength{\emergencystretch}{2em}
\multlinegap0pt
\usepackage{bm}
\usepackage{bbm}
\usepackage{dsfont}
\usepackage{xspace}
\usepackage{cancel}
\usepackage{mathtools}
\usepackage{amsthm}
\makeatletter
\@ifundefined{theorem}{\newtheorem{theorem}{Theorem}}{}
\@ifundefined{lem}{\newtheorem{lem}[theorem]{Lemma}}{}
\makeatother
\newcommand{\FD}{\mathfrak{D}}
\title[Asymmetrically Weighted Dowker Persistence and Applications ]{Asymmetrically Weighted Dowker Persistence and Applications in Dynamical Systems}
\author{Tobias Timofeyev, Christopher Potvin, Benjamin Jones, Kristin M. Kurianski, Miguel Lopez, Sunia Tanweer}
\date{Source: arXiv:2601.04559v2 (CC BY 4.0)}
\begin{document}
\maketitle

\noindent\emph{Source and licence.} Asymmetrically Weighted Dowker Persistence and Applications in Dynamical Systems. Version arXiv:2601.04559v2 (CC BY 4.0), distributed under the Creative Commons Attribution 4.0 International licence, as recorded in the arXiv licence metadata for this exact version and shown on the abstract page https://arxiv.org/abs/2601.04559v2. Full rights notice: licenses/arxiv-2601.04559-CC-BY-4.0.txt.\par\smallskip

\noindent\textit{Editorial scope.} Counted unit: the Lem whose source label is \texttt{lem:path\_wedge} in the article (source0.tex lines 1055--1060, source label \texttt{lem:path\_wedge}), delivered together with the article's own complete original proof of it (source0.tex lines 1061--1120, one \texttt{proof} environment of 60 lines). No mathematics is written, repaired or re-derived: statement and proof are the article's own text line for line. Required context retained uncounted: none. Every definition, display and label of those lines is retained unchanged. Omitted from this package and not redistributed: 1--1054, 1121--1509. Nothing inside a retained excerpt is omitted or altered: every retained body is delivered exactly as the article's lines are written, so no omission receipt of any kind accompanies this package. Citations: the retained excerpts carry no citation command at all, so no bibliography entry of the article is transcribed and no reference is redistributed. Reference adaptation: none was needed and none was made; every reference inside the delivered unit resolves inside it, so no unresolved reference or citation is produced. Printed slips kept literal: no printed word, symbol, hypothesis or number of the counted statement or of its proof was altered. Layout and class: the article's own class is \texttt{sn-jnl} with packages including lmodern, graphicx, multirow, amsmath, amssymb, amsfonts, amsthm, mathrsfs, appendix, xcolor, textcomp, manyfoot, booktabs, algorithm; the wrapper is the lane's pinned \texttt{amsart} editorial wrapper at 10pt on A4, which restates the theorem-like environments the retained text needs and adds the article's own macro definitions that the retained text needs (\texttt{\textbackslash FD}), with extra wrapper packages (bm, bbm, dsfont, xspace, cancel, mathtools). The delivered numbering is the unit's own. Assets: no retained span contains a figure environment, a graphics-inclusion command, a PDF-page inclusion, a table or a TikZ picture; no image file is copied into this package, the package declares no asset dependency and no image file is redistributed with it. Licence: version arXiv:2601.04559v2 of this article is distributed under the Creative Commons Attribution 4.0 International licence, as recorded in the arXiv licence metadata for this exact version and shown on the abstract page https://arxiv.org/abs/2601.04559v2. A full rights notice is delivered at licenses/arxiv-2601.04559-CC-BY-4.0.txt. Characters: every retained excerpt is pure ASCII, so the delivered engine, XeLaTeX with its default Latin Modern text font, reports no missing character.
\begin{lem}\label{lem:path_wedge}
    Let $G_1$ and $G_2$ be directed graphs with distinguished vertices $x_0$ and $y_0$ respectively. Let $\delta\in\mathbb{R}$. It follows that  $\tilde{H}_n(\FD^\star_\delta(P(G_1\vee G_2)))
    \approx
    \tilde{H}_n(\FD^\star_\delta(P(G_1)\vee P(G_2)))
        $.
\end{lem}

\begin{proof}
    Let $A= \text{Cl}(\FD^\star_\delta(P(G_1\vee G_2))\setminus \FD^\star_\delta(P(G_1)\vee P(G_2)))$ and $B=\FD^\star_\delta(P(G_1)\vee P(G_2))$.

    Our first claim is that $A$ is contractible. Notice that the maximal simplices of $A$ are not in $\FD^\star_\delta(P(G_1)\vee P(G_2))$. Thus they require edges between $G_1\setminus x_0$ and $G_2\setminus x_0$ in $P(G_1\vee G_2)$. Since any path in $G_1\vee G_2$ inducing such an edge must pass through $x_0$, it follows that the maximal simplices of $A$ all contain $x_0$. This induces a free face of every maximal simplex, which we may collapse with respect to. That is a maximal simplex $\sigma=[\sigma_1,\cdots,x_0,\cdots,\sigma_n]$ implies $ [\sigma_1,\cdots,\hat{x_0},\cdots,\sigma_n]$ is free. Consider the complex $A'$ where we perform a simplicial collapse with respect to this free face of every maximal simplex. It follows that the maximal simplices of $A'$ also contain $x_0$. Thus we may inductively collapse the maximal simplices until only $x_0$ remains. This shows the desired contractibility.

    Our second claim is that $A\cap B$ is contractible. This is precisely simplices returned by the closure in the definition of $A$. Let $\sigma=[\sigma_1,\dotsc,\sigma_n]$ be a maximal simplex of $A$, induced by a path from $\sigma_1\in V(G_1)\setminus x_0$ to $\sigma_n\in v(G_2)\setminus x_0$. This simplex implies $\eta_1=[\sigma_1,\dotsc,x_0]$ and $\eta_2=[x_0,\dotsc,\sigma_n]$ in $A\cap B$ representing paths to and from $x_0$, respectively. Thus all maximal simplices of $A\cap B$ contain $x_0$. Repeating the same argument as for $A$, we find that $A\cap B$  is contractible.

    Given that $\tilde{H}_n(A)=\tilde{H}_n(A\cap B)=0$, it follows by Mayer Vietoris that the following sequence is exact.

    \[
    0
    \longrightarrow
    \tilde{H}_n(B)
    \longrightarrow
    \tilde{H}_n(A\cup B)
    \longrightarrow
    0
    \]

    This implies the desired isomorphism

    \[
    \tilde{H}_n(\FD^\star_\delta(P(G_1)\vee P(G_2)))
    \approx
    \tilde{H}_n(\FD^\star_\delta(P(G_1\vee G_2))).
    \]






    % Notice that $P(G_1)\vee P(G_2)\subseteq P(G_1\vee G_2)$, and more specifically that the difference in their edge sets $E(P(G_1\vee G_2))\setminus E(P(G_1)\vee P(G_2))$ is the set of edges induced by paths from one graph to the other. 

    % Consider a single element $1-$chain $\zeta = c[x_1,x_2]$ from the chain complex on $\FD(R_\delta(P(G_1\vee G_2)))\setminus \FD(R_\delta(P(G_1)\vee P(G_2)))$ induced by a vertex $a\in V(G_1)$. We show this is homologous to a $1-$chain in $\FD(R_\delta(P(G_1)\vee P(G_2)))$. By construction $\omega(a,x_1),\omega(a,x_2)<\delta$. It follows that $\zeta$ is homologous to
    % \begin{align*}
    %     \zeta'\equiv \zeta + \partial_2(c[x_1,a,x_2])
    %     =
    %     c[x_1,a] + c[a,x_2]
    % \end{align*}
    % WLOG assume $x_1\in V(G_1)$. It follows $x_2\in V(G_2)$. If $a=x_0$ then we are done. Assume not, then it follows $[a,x_2]\in \FD(R_\delta(P(G_1\vee G_2)))\setminus \FD(R_\delta(P(G_1)\vee P(G_2)))$. Note that the path from $a\in V(G_1)$ to $x_2\in V(G_2)$ must pass through $x_0$. Thus it follows $\zeta'$ is homologous to 
    % \begin{align*}
    %     \zeta''\equiv 
    %     \zeta' + \partial_2(c[a,x_0,x_2])
    %     &=
    %     c[x_1,a] + c[x_0,x_2] + c[a,x_0]
    %     \in C_1(\FD(R_\delta(P(G_1)\vee P(G_2))))
    % \end{align*}
    % repeating this argument for all elements of a $1-$chain we see every element of $C_1(\FD(R_\delta(P(G_1\vee G_2))))$ is homologous to an element of $C_1(\FD(R_\delta(P(G_1)\vee P(G_2))))$

    % Consider an $n-$chain with $n>1$ in $C_n(\FD(R_\delta(P(G_1\vee G_2))))$ not in $C_n(\FD(R_\delta(P(G_1)\vee P(G_2))))$. Consider an $n$ simplex $\sigma$ with vertices in $G_1$ and $G_2$ that is induced by a vertex $c\in V(G_1)$. It follows that $c$ fully connects the vertices in a path form itself to a the vertices of $\sigma$ in $V(G_2)$. It follows that $\sigma$ is homologous to some $1-$chain fully contained in $\FD(R_\delta(P(G_1)))$. If this is repeated for all elements of the original $1-$chain we find it is homologous to a $1-$chain in $C_n(\FD(R_\delta(P(G_1)\vee P(G_2))))$.
    
    % The inclusion, $P(G_1)\vee P(G_2)\subseteq P(G_1\vee G_2)$, induces an inclusion map $    C_n(\FD(R_\delta(P(G_1)\vee P(G_2)))) \hookrightarrow C_n(\FD(R_\delta(P(G_1\vee G_2))))$. Consider an $n-$chain $\alpha=\sum_{i=1}^p a_i\alpha_i$ not contained in the image of this inclusion map. It follows that $a_i\neq 0$ for simplices $\alpha_i$ not in $\FD(R_\delta(P(G_1)\vee P(G_2)))$. Let $j$ be such an index. The simplex $\alpha_j$ must have been induced in the Dowker complex by an edge not in $P(G_1)\vee P(G_2)$. WLOG this edge $\omega(b_1,b_2)<\delta$ with $b_1\in V(G_1)$ and $b_2\in V(G_2)$. It follows that $\omega(a,x_0),\omega(x_0,b)<\delta$. Additionally, $[b_1,b_2,x_0]\in \FD(R_\delta(P(G_1\vee G_2)))$. Thus 

    % \begin{align*}
    %     \alpha + \partial_2([b_1,b_2,x_0])
    % \end{align*}
    
    
\end{proof}

\end{document}
